\chapter{Conclusion and Future Work}
\label{ch:conclusion}

\section{Conclusions}

We compared LoRA with full fine-tuning for ViT-B and Swin-B on CIFAR-100, in the same configurations on both models
and with five seeds.

\begin{enumerate}
  \item \textbf{Accuracy.} With the learning rate of each method chosen on validation loss, LoRA at rank 8 is more
        accurate than full fine-tuning on Swin-B (0.4--0.5~pp for each of the three targets). On ViT-B it is never
        ahead: MLP and Attn+MLP show no detectable difference, and attention LoRA is about 0.2~pp below. The comparison is sensitive to the tuning of full fine-tuning: a
        rate three times too small costs it 0.6~pp on ViT-B, enough to make LoRA look clearly better. Rank 16 gives
        no gain on ViT-B, and no target is reliably the most accurate.
  \item \textbf{Training cost.} The saving in training memory is decided by which layers are adapted, not by how
        many parameters are trained: about a third for the attention projections, a seventh for the MLP and 6\% for
        both at batch size 64, on both models, and about 70\% for any of them at batch size 1. Training epochs are
        12--26\% shorter, because the backward pass and the optimiser step become cheaper.
  \item \textbf{Placement.} Training memory is set mainly by the first adapted block; additional adapters after it
        are cheap. Adapting only the second half of the blocks saves the most memory but loses 1.2--1.6~pp and falls
        below full fine-tuning on both models. Adapting only the odd blocks, which skips the first block, saves 11--14\% of the memory with
        no detectable loss against all blocks on Swin-B; on ViT-B a small loss cannot be excluded.
  \item \textbf{Deployment.} For a single task, a merged LoRA model is indistinguishable from a fully fine-tuned
        one in predictions, memory and speed, so LoRA does not make inference cheaper. Unmerged adapters let
        several tasks share one base model at 2--6~MiB per task (adapters on all blocks), at the price of slower predictions (8--55\% for
        single images on the GPU, 2--17\% in batches of 8 or more).
  \item \textbf{Memory budgets.} Peak training memory can be predicted from the shape of the model (ViT-L within
        1.5\% before it was measured). The accuracy of a placement cannot yet: our rule, written down in advance,
        failed.
\end{enumerate}
For someone who wants to fine-tune a vision transformer with little memory, the practical summary is: put LoRA on
the attention projections, with rank 8; if memory is still short, reduce the batch size or start the adapters one or a few blocks later, not half-way through the
model; and merge the adapters before deployment unless several tasks
share the model.

\section{Limitations}
\label{sec:limitations}

\begin{itemize}
  \item \textbf{One dataset.} All results are on CIFAR-100, with images upsampled from $32\times32$. Accuracy is
        close to saturated, so differences of a few tenths of a point may not carry over to harder tasks, and the
        finding that the early blocks matter most may be specific to such images.
  \item \textbf{One task.} LoRA is usually motivated by several tasks sharing one base model. We trained one
        task, because the questions here concern the cost of adapting and deploying a single adapter. What we say
        about several tasks (stored size per task, the slower predictions of unmerged adapters) depends only on the
        shapes of the adapters and not on what they were trained for; the accuracy on other tasks is not measured.
  \item \textbf{Learning rates.} The rates were varied with three seeds and for one LoRA target. The other targets
        and all placements were trained at one rate and one rank.
  \item \textbf{Seeds and run-to-run variation.} With five seeds, differences of about 0.2~pp can be neither shown
        nor ruled out, and two sets of ViT-B runs with the same seeds differed by 0.1--0.2~pp in their means. The
        learning-rate check and the DeiT-B runs have three seeds. Merging, memory profiling and latency use one
        trained adapter per configuration.
  \item \textbf{Hardware and precision.} Everything was measured on one partition of a shared server GPU, in 32-bit
        floating point. No edge device was used; the GPU and CPU latencies show how the configurations compare, not
        what an edge device would achieve. Memory values are PyTorch's counts of memory occupied by tensors.
  \item \textbf{Power and energy were not measured.} The GPU reports power only for the whole card, which includes
        the jobs of other users on its other partitions, and the cluster does not record energy per job. In a test,
        the card drew between 280 and 503~W while our partition was idle, and a steady load on our partition raised
        the average by 6--78~W, less than this background and not repeatably.
  \item \textbf{Causes were not tested}: why LoRA loses nothing against full fine-tuning, why the best placement
        differs between the models, and why DeiT-B loses less from a late start.
\end{itemize}

\section{Future work}

\begin{itemize}
  \item \textbf{Harder and more varied tasks}, to see which accuracy results hold beyond CIFAR-100, in particular
        the importance of the early blocks.
  \item \textbf{An edge device}, including power and energy, and reduced precision, which is what such devices use.
  \item \textbf{A complete memory model and its alternatives.} The formula needs a lower bound for the temporary
        memory of the forward pass. Placement should also be compared with gradient accumulation: because memory is
        linear in the batch size, a batch can be processed in pieces with the gradients added up, which meets any
        budget with adapters on all blocks and, for ViT-B and DeiT-B, should give the same update (they normalise
        each token separately and use no dropout). We did not run this. If it holds, a later start mainly buys
        time (training epochs 14\% shorter from block 3, 28\% from block 6), so the step time of accumulation and of
        gradient checkpointing must be measured.
  \item \textbf{Predicting the accuracy of a placement}: the start-block curve at more points and on Swin-B, and a
        per-model score computed before training, such as a gradient-based sensitivity score.
\end{itemize}
